\documentclass[11pt]{scrartcl}

\usepackage[stixtwo]{fontsetup}
\usepackage{amsmath}

\begin{document}
Ha $n$ páros:
\[
	1
	\leftarrow \sqrt[n]{\frac{5}{8}} \cdot \frac{1}{\left(\sqrt[n]{n}\right)^2}
	= \sqrt[n]{\frac{5}{8n^2}}
	\le a_n
	= \sqrt[n]{\frac{5}{2n^2 + 6}}
	\le \sqrt[n]{\frac{5}{n^2}}
	= \frac{\sqrt[n]{5}}{\left(\sqrt[n]{n}\right)^2}
	\to 1,
\]
ha pedig $n$ páratlan, akkor
\[
	1
	\leftarrow \sqrt[n]{\frac{2}{6}} \left(\sqrt[n]{n}\right)^3
	= \sqrt[n]{\frac{2n^3}{6}}
	\le a_n
	= \sqrt[n]{\frac{2n^3 + 5}{6}}
	\le \sqrt[n]{\frac{7n^3}{6}}
	= \sqrt[n]{\frac{7}{6}} \left(\sqrt[n]{n}\right)^3
	\to 1,
\]
így a sorozat egyetlen torlódási pontja $1$, tehát
$\limsup a_n = 1 = \liminf a_n = \lim a_n$.

Páros $n$ esetén
\[
	0 \le b_n = \sqrt{\frac{5}{2n^2 + 6}} \le \sqrt{\frac{5}{2}} \cdot \frac{1}{n} \to 0,
\]
páratlan $n$ esetén
\[
	b_n = \sqrt{\frac{2n^3 + 5}{6}} \ge \sqrt{\frac{1}{3}} \cdot n^{\frac{3}{2}} \to \infty,
\]
tehát a két torlódási pont $0$ és $\infty$, így $\limsup b_n = \infty$,
$\liminf b_n = 0$, és a sorozat divergens.
\end{document}
